% !TeX root = ../../../../../main.tex
% chapters/sections/chapter5/subsections/asad/as.tex

\subsection{総供給（AS）曲線}
\label{sec:chapter5_asad_as}
供給側は自国代表的家計の価格の最適化を記述する以下の4本の非線形方程式によって構成される。
\begin{itemize}
    \item 最適価格決定式\eqref{form:chapter3_representative_household_phillips_home_representative_p_H_tilde_def}
    \item 補助変数\(v_t\)の再帰動学式\eqref{form:chapter3_representative_household_phillips_home_representative_v_H_eq}
    \item 補助変数\(w_t\)の再帰動学式\eqref{form:chapter3_representative_household_phillips_home_representative_w_H_eq}
    \item 正規化された生産者物価指数（PPI）\(\bar{p}_t^H\)の動学方程式\eqref{form:chapter3_representative_household_phillips_home_representative_p_H_bar_dynamics_p_H_bar_modified2_eq}
\end{itemize}
これらの方程式を定常状態の周りで対数線形近似し補助変数を消去して整理することにより、
以下の線形化されたフィリップス曲線が得られる（付録\ref{TODO}）。
\begin{equation}
\label{form:chapter5_asad_as_pi_H_eq}
\hat{\pi}_t^H=\beta_{ss}^HE_t[\hat{\pi}_{t+1}^H]+\kappa^H(\hat{y}_t^H-2\hat{a}_t^H-\hat{p}_t^H-\hat{\lambda}_t^H-\hat{\tau}_t^H)
\end{equation}
ここでハット記号\(\hat{}\)は定常状態からの対数乖離率を、下付き文字の\(ss\)は定常状態の値を意味する。
\(\kappa^H\)は式\eqref{appendixTODO}で定義される価格の粘着性に依存する正の定数である。
\(\hat{\pi}_t^H\)は式\eqref{appendixTODO}で以下のように定義されるインフレ率である。
\begin{equation}
\label{form:chapter5_asad_as_pi_H_def}
    \hat{\pi}_t^H\equiv\hat{p}_t^H-\hat{p}_{t-1}^H
\end{equation}
パラメータ設定の表\ref{tab:chapter4_parameters_list}でも述べたとおり、
本稿では日本経済の実証データにもとづき自国の価格は極めて硬直的に設定している（\(\xi^H=0.99\)）。
この条件下では\(\kappa^H\)が極めて小さいため本稿の分析では単純化して\(\kappa^H=0\)とおく。
このときインフレ率\(\hat{\pi}_t^H\)は期待インフレ率\(E_t[\hat{\pi}_{t+1}^H]\)にのみ依存する形式となる。
ここで定義式\eqref{form:chapter5_asad_as_pi_def}より
\(\hat{\pi}_t^H=\hat{p}_t^H-\hat{p}_{t-1}^H\)および
\(E_t[\hat{\pi}_{t+1}^H]=E_t[\hat{p}_{t+1}^H]-\hat{p}_t^H\)
となるから、これらを代入して\(\hat{p}_t^H\)について整理すると以下の\textbf{AS曲線}が得られる。
\begin{equation}
\label{form:chapter5_asad_as_p_H_eq}
\hat{p}_t^H=\frac{1}{1+\beta_{ss}^H}\hat{p}_{t-1}^H+\frac{\beta_{ss}^H}{1+\beta_{ss}^H}E_t[\hat{p}_{t+1}^H]
\end{equation}
この式には生産\(\hat{y}_t^H\)が含まれていない。
したがって\((y,p)\)平面においてAS曲線は\textbf{水平な直線}として描かれる。
この水平線の高さは前期の生産者物価指数\(\hat{p}_{t-1}^H\)と
期待生産者物価指数\(E_t[\hat{p}_{t+1}^H]\)のみによって決定される。